\begin{python0}
from solutions import *; clear()
\end{python0}
\sectionthree{Example: counting digits}

Here's a problem where the condition does not depend on
continually prompting the user for an input. Therefore there is
\EMPHASIZE{no sentinel value}. You will see that there is also no control
variable -- so you definitely should not expect a for-loop.

Pay attention to the way I slowly work through the problem until I see
the loop.

Suppose the user enters an integer and you want to count the number of
digits in that integer.

For instance if the user enters 134, your program must display 3. If the
user enters 97531, your program must display 5.

So how do you do that? If you count the digits from the right to left,
you can think of moving a finger one digit at a time and counting each
digit your finger passed by:

%ISSUE: whatever the best way of actually displaying ^ is can be found later when less important things need to be done
\tab[3em]{97531\^{}}\\
\tab[3em]{9753\^{}$^{1}$\tab[3em]{1 digit}}\\
\tab[3em]{975\^{}$^{31}$ \tab[3em]{2 digits}}\\
\tab[3em]{97\^{}$^{531}$   \tab[3em]{3 digits}}\\
\tab[3em]{9\^{}$^{7531}$     \tab[3em]{4 digits}}\\
\tab[3em]{\^{}$^{97531}$      \tab[3em]{5 digits}}\\
\tab[3em]{(\scalebox{2}{$^{}$} = your finger.)}\\

Now you notice that the digits you've passed does not
affect the rest of your digit-counting process. So let's
throw them away:\\
\tab[3em]{97531\^{}}\\
\tab[3em]{9753\^{}\tab[3em]{1 digit}}\\
\tab[3em]{975\^{} \tab[3em]{2 digits}}\\
\tab[3em]{97\^{}   \tab[3em]{3 digits}}\\
\tab[3em]{9\^{}     \tab[3em]{4 digits}}\\
\tab[3em]{\^{}      \tab[3em]{5 digits}}\\

Now, the symbol for our finger doesn't help.
Let's get rid of that too:\\

\begin{align*}
97531\\
9753 &\quad\text{1 digit}\\
975 &\quad\text{2 digits}\\
97 &\quad\text{3 digits}\\
9 &\quad\text{4 digits}\\
\end{align*}
Well ... that sure looks like removing the rightmost digit (one at a
time) to me. And we already know how to do that: just do integer
division by 10!!! If we do that, then the last number should be zero:\\
\begin{align*}
97531\\
9753 &\quad\text{1 digit}\\
975 &\quad\text{2 digits}\\
97 &\quad\text{3 digits}\\
9 &\quad\text{4 digits}\\
0 &\quad\text{5 digits}\\
\end{align*}
Of course the number that is going up (the 1, 2, 3, 4, 5) requires a
variable. Let's call it x. And let's say
the number we're working on (the number entered by the
user) is placed in variable n.

Now when does the process \EMPHASIZE{stop}? When we see a zero in the left
column of the above simulation, right? So the process seems to be this:
\begin{align*}
\text{prompt user for integer n}\\
\text{if n is not zero, n = n / 10 and x = 1}\\
\text{if n is not zero, n = n / 10 and x = 2}\\
\text{if n is not zero, n = n / 10 and x = 3}\\
\text{if n is not zero, n = n / 10 and x = 4}\\
\text{if n is not zero, n = n / 10 and x = 5}\\
\text{n is zero. Stop.}\\
\end{align*}
We need to write this as an algorithm. Obviously there's
a loop here!!! You want the statements
\begin{align*}
\text{if n is not zero, n = n / 10 and x = 1}\\
\text{if n is not zero, n = n / 10 and x = 2}\\
\text{if n is not zero, n = n / 10 and x = 3}\\
\text{if n is not zero, n = n / 10 and x = 4}\\
\text{if n is not zero, n = n / 10 and x = 5}\\
\end{align*}
to be replaced by one statement. You can ``combine'' them like this:

\begin{align*}
\text{if n is not zero, n = n / 10 and x = x + 1}\\
\text{if n is not zero, n = n / 10 and x = x + 1}\\
\text{if n is not zero, n = n / 10 and x = x + 1}\\
\text{if n is not zero, n = n / 10 and x = x + 1}\\
\text{if n is not zero, n = n / 10 and x = x + 1}\\
\end{align*}

AHA! They are all the same! Of course x does not have an initial value.
You need to initialize it. And based on the above example, you can see
that x should be initialized to 0. Now you can write the following very
easily:
\begin{align*}
\text{while n is not 0:}\\
&\quad\text{n = n / 10}\\
&\quad\text{x = x + 1}\\
\end{align*}
Here's the code:

\begin{console}
int n = 97531;
int numDigits = 0;

while (n != 0)
{     
      n /= 10;
      ++numDigits;
      std::cout << n << ' '<< numDigits << std::endl;
}

std::cout << numDigits << std::endl; 
\end{console}

I've inserted print statements to get the program to show us the intermediate steps. And instead of \texttt{x}, I'm using \texttt{numDigits} which is a lot more descriptive.

Yes, of course you can rewrite the while-loop as a for-loop:

\begin{console}
int numDigits;

for (numDigits = 0; n != 0; ++numDigits)
{   
    n /= 10;
    std::cout << n << ' ' << numDigits << std::endl;
} 
\end{console}

Although it's not wrong, note that the for-loop does not really use a control variable: the variable updated (i.e. numDigits) is the variable in the boolean expression (i.e. n). This is an indication that the while-loop is probably best.

Remember: If the loop doesn't look like a for-loop
(i.e., there's no single variable that runs across a
sequence of values), then don't use it. Use the
while-loop instead.

